\subsection{Cel i zakres pracy}

Celem pracy było porównanie wydajności czterech systemów zarządzania
bazami danych -- PostgreSQL, MySQL, MongoDB i Neo4j -- na realistycznym
modelu danych platformy streamingowej VOD. Porównanie przeprowadzono
przez 24 scenariusze CRUD (po 6 z każdej kategorii: INSERT, SELECT,
UPDATE, DELETE), wykonywane na trzech wolumenach danych
(\emph{small} -- 500\,tys., \emph{medium} -- 1\,mln, \emph{large} --
10\,mln rekordów), w dwóch wariantach indeksowych (bez indeksów
dodatkowych i z nimi). Centralnym pytaniem badawczym była weryfikacja
hipotezy o wpływie indeksów na operacje czytające i zapisujące.

% -------------------------------------------------------
\subsection{Co zostało wykonane}

W ramach pracy zaprojektowano cztery równoważne reprezentacje tej
samej domeny -- relacyjną dla PostgreSQL i MySQL, dokumentową dla
MongoDB oraz grafową dla Neo4j -- co opisano w rozdziale 6. Stworzono
deterministyczny generator danych zapisujący wyniki do plików CSV
zaczerpywanych przez cztery niezależne loadery (rozdział 8). Zbudowano
warstwę orkiestracji w Pythonie (rozdziały 7 i 9) odpowiedzialną za
sekwencyjne uruchamianie pipeline'u: generowanie~$\rightarrow$~ładowanie~$\rightarrow$~benchmark~$\rightarrow$
analiza EXPLAIN~$\rightarrow$~wizualizacja, w dwóch wariantach indeksowych.
Łącznie wykonano \textbf{1\,728 pomiarów czasów wykonania} oraz
\textbf{144 analizy planów wykonania} (po 6 scenariuszy SELECT
$\times$~4 silniki $\times$~3 wolumeny $\times$~2 warianty), z których
zsyntetyzowano analizę przedstawioną w rozdziałach 11--13.

% -------------------------------------------------------
\subsection{Główne wnioski}

\textbf{Hipoteza została potwierdzona w postaci ogólnej.} Indeksy
istotnie przyspieszają operacje SELECT (najmocniejsze redukcje rzędu
$10^2$--$10^3$ razy w scenariuszach S5 MongoDB i D6 PostgreSQL) oraz
spowalniają masowe operacje UPDATE i INSERT na kolumnach indeksowanych
(scenariusz U3 PostgreSQL: 2.4-krotne spowolnienie). Skala efektu
zależy jednak silnie od dwóch czynników -- selektywności filtru oraz
liczby modyfikowanych wierszy w kolumnach indeksowanych. Pełna
weryfikacja punkt po punkcie znajduje się w rozdziale 13.

\textbf{Architektura silnika ma porównywalne znaczenie z obecnością
indeksów.} W scenariuszu S2 (collaborative filtering) Neo4j na
wolumenie \emph{large} jest 8--15-krotnie szybszy od pozostałych baz,
niezależnie od wariantu indeksowego -- bo wielokrotna nawigacja
po krawędziach jest tańsza od kolejnych \texttt{JOIN}-ów po milionach
wierszy. Z drugiej strony MongoDB w scenariuszach \texttt{INSERT}
masowych (I2) demonstruje ograniczenia modelu dokumentowego -- 3-krotne
spowolnienie po założeniu indeksu. Każdy z czterech paradygmatów ma
,,słodkie punkty'' i miejsca, w których strukturalnie traci.

\textbf{Operacje DELETE pokazują, że hipoteza ma postać ogólniejszą.}
Mimo że pierwotna hipoteza nie obejmowała operacji usuwających, dane
wskazują, że indeks pomaga w każdej operacji wyposażonej w \emph{fazę
wyszukiwania} -- włącznie z DELETE. Najsilniejszy wynik całej pracy
(D6 PostgreSQL, redukcja z 30.7~sekundy do 10~milisekund) dotyczy
właśnie tej kategorii. Uogólniona forma hipotezy: \emph{indeksy
przyspieszają fazę wyszukiwania w dowolnej operacji CRUD i spowalniają
fazę modyfikacji w operacjach modyfikujących -- przewaga albo
penalizacja zależy od proporcji między tymi fazami}.

\textbf{Praktyczna rekomendacja}: indeksy należy
zakładać agresywnie dla kolumn intensywnie filtrowanych w SELECT
i DELETE, ostrożnie dla kolumn często aktualizowanych masowo;
najwyższy priorytet powinny mieć indeksy na kluczach obcych w tabelach
biorących udział w usuwaniu kaskadowym. Wybór silnika powinien być
podyktowany dominującym wzorcem dostępu w aplikacji -- dla wzorców
nawigacyjnych (rekomendacje, sieci powiązań) wygrywa graf, dla
agregacji punktowych wygrywa baza relacyjna z bogatym zestawem
indeksów (PostgreSQL), dla pojedynczych operacji \texttt{INSERT}
struktur zagnieżdżonych wygrywa model dokumentowy.

% -------------------------------------------------------
\subsection{Ograniczenia pracy}

Pomiary były wykonywane jednowątkowo, bez współbieżnego obciążenia,
co nie odzwierciedla typowego scenariusza produkcyjnego, w którym
silnik obsługuje wiele równoczesnych klientów. Liczba prób na
scenariusz (\texttt{trials}=3) jest mała statystycznie -- mediana
z trzech wartości jest odporna na pojedyncze odchylenia, ale nie
pozwala na wnioski o rozkładzie ogonów. Konfiguracja silników została
świadomie zbliżona do siebie pod względem limitów pamięci, ale nie
była dostrojona per silnik -- w realnym wdrożeniu każdy z nich
zasługiwałby na własną, dedykowaną kalibrację. MySQL nie pozwala
upuścić indeksu chroniącego klucz obcy, przez co baseline
,,bez indeksów'' jest dla niego częściowo skażony i utrudnia
bezpośrednie porównanie z pozostałymi silnikami. Wreszcie, 24
scenariusze pokrywają reprezentatywny przekrój operacji platformy
VOD, ale nie wyczerpują wszystkich możliwych wzorców dostępu --
w szczególności brakuje zapytań analitycznych typu OLAP, agregacji
po długich oknach czasowych oraz operacji geoprzestrzennych.

% -------------------------------------------------------
